\section*{Capítulo \ref{ch:g12:meiosis-genetic-shuffling} --- A meiose e o embaralhamento genético}

\begin{solution}{exo:g12:meiosis-genetic-shuffling:1}
\omterm{def:g12:meiosis-genetic-shuffling:meiosis}{Diploide}: duas cópias de cada \omterm{def:g11:cell-cycle-mitosis:chromosome}{cromossomo} ($2n$); \omterm{def:g12:meiosis-genetic-shuffling:meiosis}{haploide}: uma ($n$). A
\omterm{def:g12:meiosis-genetic-shuffling:meiosis}{meiose} divide pela metade o número de \omterm{def:g10:universal-dna:information}{cromossomos}, de $2n$ para $n$, e
reduz o \omterm{def:g10:universal-dna:information}{DNA} de $2q$ (depois da replicação) para $q/2$ por gameta.
\end{solution}

\begin{solution}{exo:g12:meiosis-genetic-shuffling:2}
A \omterm{def:g12:meiosis-genetic-shuffling:meiosis}{meiose} I separa os \omterm{def:g10:universal-dna:information}{cromossomos} homólogos de cada par; a \omterm{def:g12:meiosis-genetic-shuffling:meiosis}{meiose} II
separa as duas \omterm{def:g11:cell-cycle-mitosis:chromosome}{cromátides} de cada \omterm{def:g11:cell-cycle-mitosis:chromosome}{cromossomo}. A primeira é a reducional.
\end{solution}

\begin{solution}{exo:g12:meiosis-genetic-shuffling:3}
Depois da \omterm{def:g12:meiosis-genetic-shuffling:meiosis}{meiose} I: 4 \omterm{def:g10:universal-dna:information}{cromossomos} (duas \omterm{def:g11:cell-cycle-mitosis:chromosome}{cromátides} cada), \omterm{def:g10:universal-dna:information}{DNA} $q$.
Depois da \omterm{def:g12:meiosis-genetic-shuffling:meiosis}{meiose} II: 4 \omterm{def:g10:universal-dna:information}{cromossomos} (uma \omterm{def:g11:cell-cycle-mitosis:chromosome}{cromátide} cada), \omterm{def:g10:universal-dna:information}{DNA} $q/2$.
\end{solution}

\begin{solution}{exo:g12:meiosis-genetic-shuffling:4}
$2^4 = 16$.
\end{solution}

\begin{solution}{exo:g12:meiosis-genetic-shuffling:5}
Um cruzamento com um parceiro que carrega apenas \omterm{def:g10:universal-dna:gene}{alelos} \omterm{def:g11:genes-and-disease:genetic}{recessivos}, de
modo que cada descendente mostra o gameta do genitor testado. A
proporção de recombinantes revela se dois \omterm{def:g10:universal-dna:gene}{genes} estão no mesmo
\omterm{def:g11:cell-cycle-mitosis:chromosome}{cromossomo} e, em caso afirmativo, a que distância um do outro.
\end{solution}

\begin{solution}{exo:g12:meiosis-genetic-shuffling:6}
\omterm{def:g12:meiosis-genetic-shuffling:meiosis}{Meiose} I: $2q$; entre as divisões: $q$; gameta: $q/2$. Sem uma fase S
intercalada, a segunda divisão divide o \omterm{def:g10:universal-dna:information}{DNA} pela metade de novo, e é
isso que leva o gameta à metade do conteúdo de uma \omterm{def:g10:cells-common-unit:cell}{célula} \omterm{def:g12:meiosis-genetic-shuffling:meiosis}{diploide}.
\end{solution}

\begin{solution}{exo:g12:meiosis-genetic-shuffling:7}
Quatro classes iguais: não estão ligados, estão em \omterm{def:g10:universal-dna:information}{cromossomos}
diferentes. Gametas $AB$, $Ab$, $aB$, $ab$, um quarto cada.
\end{solution}

\begin{solution}{exo:g12:meiosis-genetic-shuffling:8}
Ligados: duas classes grandes ($AB$, $ab$, parentais) e duas pequenas
($Ab$, $aB$, recombinantes). Recombinantes $85/1000 = 8.5\%$.
\end{solution}

\begin{solution}{exo:g12:meiosis-genetic-shuffling:9}
Apenas $AB$ e $ab$: uma troca acima dos dois \omterm{def:g10:universal-dna:gene}{genes} troca segmentos que
não carregam nenhum dos dois, de modo que as combinações de $A/a$ e de
$B/b$ não mudam. A recombinação entre dois \omterm{def:g10:universal-dna:gene}{genes} exige uma troca entre
eles.
\end{solution}

\begin{solution}{exo:g12:meiosis-genetic-shuffling:10}
Cerca de 0.8, 2.9 e 16 por 1000: um fator de 20 entre 25 e 42 anos.
\end{solution}

\begin{solution}{exo:g12:meiosis-genetic-shuffling:11}
Na \omterm{def:g12:meiosis-genetic-shuffling:meiosis}{meiose} I, os dois homólogos vão para uma mesma \omterm{def:g10:cells-common-unit:cell}{célula}: dois gametas
com duas cópias e dois sem nenhuma. Na \omterm{def:g12:meiosis-genetic-shuffling:meiosis}{meiose} II, as duas \omterm{def:g11:cell-cycle-mitosis:chromosome}{cromátides} de
um \omterm{def:g11:cell-cycle-mitosis:chromosome}{cromossomo} vão para uma mesma \omterm{def:g10:cells-common-unit:cell}{célula}: um gameta com duas cópias, um
sem nenhuma, e dois normais.
\end{solution}

\begin{solution}{exo:g12:meiosis-genetic-shuffling:12}
A recombinação é proporcional à distância: 2\% quer dizer que os \omterm{def:g10:universal-dna:gene}{genes}
estão muito próximos, de modo que raramente uma troca cai entre eles;
48\% quer dizer bem afastados, com uma troca entre eles em quase toda
\omterm{def:g12:meiosis-genetic-shuffling:meiosis}{meiose}. Como uma troca dá 50\% de gametas recombinantes, \omterm{def:g10:universal-dna:gene}{genes} ligados
e distantes recombinam com a mesma liberdade que \omterm{def:g10:universal-dna:gene}{genes} independentes.
\end{solution}

\begin{solution}{exo:g12:meiosis-genetic-shuffling:13}
$AB$, $Ab$, $aB$, $ab$, um quarto cada, pela segregação independente
dos dois pares. A \omterm{prop:g12:meiosis-genetic-shuffling:crossingover}{permutação} não muda nada para esses \omterm{def:g10:universal-dna:gene}{genes}: ela
recombina \omterm{def:g10:universal-dna:gene}{alelos} ao longo de um mesmo \omterm{def:g11:cell-cycle-mitosis:chromosome}{cromossomo}, e estes estão em
\omterm{def:g10:universal-dna:information}{cromossomos} diferentes.
\end{solution}

\begin{solution}{exo:g12:meiosis-genetic-shuffling:14}
Cada genitor passa a cada filho um de dois \omterm{def:g10:universal-dna:gene}{alelos} em cada \omterm{def:g10:universal-dna:gene}{gene},
escolhido ao acaso; num dado \omterm{def:g10:universal-dna:gene}{gene}, dois irmãos receberam o mesmo \omterm{def:g10:universal-dna:gene}{alelo}
de um genitor com probabilidade $1/2$. Ao longo de muitos \omterm{def:g10:universal-dna:gene}{genes} a
fração compartilhada tem média $1/2$, mas para um par específico de
irmãos ela varia em torno disso.
\end{solution}

\begin{solution}{exo:g12:meiosis-genetic-shuffling:15}
Clones da planta-mãe, geneticamente idênticos a ela e entre si. A
planta ganha uma multiplicação rápida e segura de um \omterm{def:g11:enzymes-and-phenotype:phenotype}{genótipo} que
funciona aqui e agora; ela perde a variedade que permitiria a alguns
descendentes sobreviver a uma mudança de ambiente ou a um parasita novo.
\end{solution}

\begin{solution}{pb:g12:meiosis-genetic-shuffling:1}
\textbf{1.} Fêmea $b^+b$, $vg^+vg$, com $b^+vg^+$ num \omterm{def:g11:cell-cycle-mitosis:chromosome}{cromossomo} e
$b\,vg$ no outro (vindos dos dois genitores puros dela). Macho $bb$,
$vg\,vg$.

\textbf{2.} Apenas gametas $b\,vg$: o macho contribui só com \omterm{def:g10:universal-dna:gene}{alelos}
\omterm{def:g11:genes-and-disease:genetic}{recessivos}, de modo que o \omterm{def:g11:enzymes-and-phenotype:phenotype}{fenótipo} de cada descendente revela o gameta
da fêmea.

\textbf{3.} $b^+vg^+$ (cinza-normal), $b\,vg$ (preto-vestigial),
$b^+vg$ (cinza-vestigial), $b\,vg^+$ (preto-normal).

\textbf{4.} Cinza-normais 42\%, pretos-vestigiais 41\%,
cinza-vestigiais 9\%, pretos-normais 8\%.

\textbf{5.} No mesmo \omterm{def:g11:cell-cycle-mitosis:chromosome}{cromossomo}: duas classes grandes e duas pequenas,
em vez de $1:1:1:1$.

\textbf{6.} Cinza-normal e preto-vestigial: as combinações carregadas
pelos dois homólogos dela, herdadas intactas dos genitores puros dela.

\textbf{7.} Cinza-vestigial e preto-normal, vindas de uma \omterm{prop:g12:meiosis-genetic-shuffling:crossingover}{permutação}
entre os dois \omterm{def:g10:universal-dna:gene}{genes} na prófase I.

\textbf{8.} $(206 + 185)/2300 \approx 17\%$. Uma \omterm{prop:g12:meiosis-genetic-shuffling:crossingover}{permutação} entre os
\omterm{def:g10:universal-dna:gene}{genes} torna duas das quatro \omterm{def:g11:cell-cycle-mitosis:chromosome}{cromátides} recombinantes, de modo que 17\%
de gametas recombinantes querem dizer uma troca em cerca de 34\% das
\omterm{def:g12:meiosis-genetic-shuffling:meiosis}{meioses}.

\textbf{9.} Uma única troca produz as duas \omterm{def:g11:cell-cycle-mitosis:chromosome}{cromátides} recombinantes
juntas, uma de cada tipo; e os dois homólogos vão para os gametas por
igual, de modo que as duas classes parentais também se equivalem.

\textbf{10.} As classes parentais passam a ser cinza-vestigial e
preto-normal, cerca de 41.5\% cada; as recombinantes, cinza-normal e
preto-vestigial, cerca de 8.5\% cada.

\textbf{11.} Sim: recombinação bem abaixo de 50\% com os dois.

\textbf{12.} $b$ --- $cn$ --- $vg$: 9 e 8 somam 17.

\textbf{13.} A chance de uma troca num intervalo é proporcional ao
comprimento dele, de modo que as frequências de intervalos adjacentes
se somam: a frequência de recombinação mede a distância ao longo do
\omterm{def:g11:cell-cycle-mitosis:chromosome}{cromossomo}.

\textbf{14.} Uma troca dá duas \omterm{def:g11:cell-cycle-mitosis:chromosome}{cromátides} recombinantes em quatro: no
máximo 50\%; trocas adicionais reembaralham, mas nunca elevam a fração
recombinante acima da metade.

\textbf{15.} Ele está em outro \omterm{def:g11:cell-cycle-mitosis:chromosome}{cromossomo} ou muito longe no mesmo; os
cruzamentos não conseguem dizer qual das duas coisas.

\textbf{16.} $2^4 = 16$ gametas; $16 \times 16 = 256$ combinações.

\textbf{17.} Toda troca produz \omterm{def:g10:universal-dna:information}{cromossomos} que nunca existiram antes,
numa posição que varia de uma \omterm{def:g12:meiosis-genetic-shuffling:meiosis}{meiose} para a seguinte; o número de
\omterm{def:g10:universal-dna:information}{cromossomos} distintos, e portanto de gametas, não tem limite prático.

\textbf{18.} $2^{23} \times 2^{23} \approx \num{7e13}$: setecentas
vezes mais combinações que o número de seres humanos que já viveram, e
isso antes da \omterm{prop:g12:meiosis-genetic-shuffling:crossingover}{permutação}.

\textbf{19.} Os irmãos tiram os \omterm{def:g10:universal-dna:gene}{alelos} deles dos mesmos quatro
conjuntos (dois por genitor) e compartilham, em média, metade deles; os
estranhos tiram de conjuntos diferentes.

\textbf{20.} Os 17\%, a fração dos gametas da fêmea recombinados entre
$b$ e $vg$, medem a distância entre os \omterm{def:g10:universal-dna:gene}{genes}; a segregação independente
dos pares e a \omterm{prop:g12:meiosis-genetic-shuffling:crossingover}{permutação} dentro de cada par fazem cada gameta ser
diferente.
\end{solution}
